\documentclass[12pt]{article}
\usepackage{amsmath, amssymb}
\usepackage{geometry}
\usepackage{enumitem}
\geometry{margin=1in}

\title{Trigonometry Mixed Review (Algebra II)}
\author{}
\date{}

\begin{document}
\maketitle

\section*{Instructions}
Complete in 60 minutes. Show clear work. Use exact values when possible; otherwise round to the nearest tenth.

\section*{Reference Facts}
\begin{itemize}[leftmargin=1.5em]
  \item Triangle area: $K=\dfrac{1}{2}ab\sin C$
  \item Law of Sines: $\dfrac{a}{\sin A}=\dfrac{b}{\sin B}=\dfrac{c}{\sin C}$
  \item Law of Cosines: $c^2=a^2+b^2-2ab\cos C$
  \item Principal inverse trig ranges:
  \[
  \arcsin x\in\left[-\frac{\pi}{2},\frac{\pi}{2}\right],\quad
  \arccos x\in[0,\pi],\quad
  \arctan x\in\left(-\frac{\pi}{2},\frac{\pi}{2}\right)
  \]
\end{itemize}

\section*{Problems}
\begin{enumerate}[leftmargin=*, itemsep=1em, topsep=0.8em]
  \item Find the exact value of
  \[
  \frac{\sin\left(\frac{7\pi}{6}\right)\cos\left(\frac{5\pi}{4}\right)}{\tan\left(\frac{5\pi}{3}\right)}
  +\sec\left(\frac{2\pi}{3}\right).
  \]

  \item Find the exact value of
  \[
  \frac{\csc\left(\frac{11\pi}{6}\right)-\cot\left(\frac{5\pi}{6}\right)}{\sec\left(\frac{5\pi}{4}\right)}.
  \]

  \item If $\sin\theta=\dfrac{3}{5}$ and $\cos\phi=-\dfrac{12}{13}$ where $\theta$ is in Quadrant II and $\phi$ is in Quadrant III, find the exact values of
  \[
  \cos\theta,\qquad \tan\phi,\qquad \sin(\theta+\phi).
  \]

  \item Find the exact values:
  \[
  \arcsin\left(-\frac{1}{2}\right),\qquad
  \arccos\left(-\frac{\sqrt{2}}{2}\right),\qquad
  \arctan(-\sqrt{3}),
  \]
  then compute
  \[
  \arccos\left(-\frac{\sqrt{2}}{2}\right)-\arcsin\left(-\frac{1}{2}\right).
  \]

  \item Evaluate exactly:
  \[
  \sin\!\left(2\arccos\frac{5}{13}\right),\qquad
  \tan\!\left(\arcsin\frac{8}{17}\right),\qquad
  \cos\!\left(\arctan\left(-\frac{3}{4}\right)\right).
  \]

  \item Two sides of a triangle are 19 cm and 27 cm, and the included angle is $132^\circ$. Find the area of the triangle and the length of the third side.

  \item In $\triangle ABC$, $a=17$, $b=24$, and $C=41^\circ$. Find $c$, then find the area of the triangle, then determine the larger of $\angle A$ and $\angle B$ to the nearest tenth of a degree.

  \item In $\triangle PQR$, $p=14$, $q=20$, and $P=32^\circ$. Determine how many triangles are possible. Solve every possible triangle.

  \item In $\triangle XYZ$, $x=11$, $y=15$, and $z=23$. Solve the triangle and classify it as acute, right, or obtuse.

  \item For
  \[
  y=-4\sin\left(3x+\frac{2\pi}{3}\right)+2,
  \]
  state the amplitude, period, phase shift, midline, and range. Then give one full period interval and the $x$-coordinates of one maximum and one minimum on that interval.

  \item For
  \[
  y=3\cos\left(\frac{x-\pi}{2}\right)-1,
  \]
  state the amplitude, period, phase shift, midline, and range. Then find all $x$ in $[0,4\pi]$ for which $y=2$.

  \item For each function below, state the period, all vertical asymptotes in one period, and all zeros in one period:
  \[
  y=\tan\left(2x-\frac{\pi}{3}\right),\qquad
  y=-\cot\left(x+\frac{\pi}{4}\right),\qquad
  y=2\sec\left(\frac{x}{2}\right)+1,\qquad
  y=\csc\left(3x-\frac{\pi}{2}\right).
  \]
\end{enumerate}

\end{document}
